\subsection{C.3 over every finite commutative coefficient ring}
\label{C:finite-bases-subsection}

Here a \emph{finite ring} has a finite underlying set, supplied
by its addition and multiplication tables. The construction uses finite
sets of polynomial classes, with no freeness assertion about a bounded
coefficient coalgebra over the base.

\begin{lemma}[Known residual finiteness, with an explicit equality algorithm]
\label{C:finite-base-polynomial-equality}
Let $R$ be a finite commutative ring, and let
$A=R[z_1,\ldots,z_s]/(f_1,\ldots,f_t)$ be a specified finite presentation.
Every nonzero element of $A$ survives in a finite quotient ring.
Equality of two polynomial classes in $A$ is decidable by a uniform
terminating algorithm from the displayed finite data. The assertion
also applies to tensor products of two such finitely presented
$R$-algebras.
\end{lemma}
\begin{proof}
We first give the residual-finiteness argument, also recorded for all
finitely generated commutative rings in
\cite[\S2.1, before Lemma 2.1]{C-JaikinZapirain2023}.
The ring $A$ is Noetherian. Given $0\ne a\in A$, choose a maximal ideal
$\mathfrak m$ containing $\operatorname{Ann}(a)$. Then $a/1\ne0$ in
$A_{\mathfrak m}$: an element outside $\mathfrak m$ cannot annihilate
$a$. Krull intersection gives $a/1\notin\mathfrak m^N A_{\mathfrak m}$
for some $N$, so $a\notin\mathfrak m^N$
\cite[Tag 00IP, Lemma 10.51.4]{C-StacksFiniteBaseEquality}.
The image of $R$ in $A/\mathfrak m$ is a finite domain, hence a finite
field. The latter residue field is a finite-type algebra over that
field, so it is finite by the Nullstellensatz
\cite[Tag 00FV, Theorem 10.34.1(1)]{C-StacksFiniteBaseEquality}.
Every $\mathfrak m^j/\mathfrak m^{j+1}$ is a finite-dimensional vector
space over $A/\mathfrak m$, since $\mathfrak m$ is finitely generated.
It follows that $A/\mathfrak m^N$ is finite. This quotient separates $a$.

For the algorithm, subtract the two input polynomials and call the
result $h$. At stage $b=0,1,2,\ldots$ perform both of the following
finite searches.
\begin{enumerate}
\item Enumerate all $t$-tuples of polynomials $g_i$ of degree at most
$b$ over $R$, and check the formal polynomial equality
$h=\sum_i g_i f_i$. If found, report equality in $A$.
\item Enumerate all finite commutative unital ring tables with at most
$b$ elements, all unital ring maps from $R$ to those rings, and all
$s$-tuples of elements in them. Retain evaluations for which every
$f_i$ is zero. If one has $h\ne0$, report inequality in $A$.
\end{enumerate}
All ring axioms, homomorphism conditions and polynomial evaluations
are finite table checks. The two reports cannot conflict. If $h=0$
in $A$, its membership in the finitely generated defining ideal has
a finite polynomial witness, so the first search eventually reports.
If $h\ne0$, the finite quotient just proved supplies an evaluation
eventually found by the second search. Thus one search terminates,
with a correct answer, on every input. This is a decision algorithm,
not merely a semidecision procedure. Finally, a tensor product has the
finite presentation with both sets of variables and defining
relations, so the same argument applies.
\end{proof}

\begin{theorem}[Proved: complete finite-base indices, additive groups and operations]
\label{C:Om-finite-base-algorithm}
For every finite commutative ring $R$ with $2\in R^\times$ and every
split $O_m$ over $R$, there is a uniform algorithm giving explicit
canonical indices for all indecomposable finite projective rational
representations, all their duality and residue-form data, and the
entire ring $WG_{\mathrm{sp}}(BO_m/R)$ with its exterior operations.

More precisely, write the computable finite local decomposition as
$R=\prod_{\nu=1}^sR_\nu$. Over each $R_\nu$, the index set
$\mathscr J(R_\nu)$ is a union of explicit finite lists, one for each
free rank $n\geq1$ and minimal polynomial coefficient degree $d\geq0$.
For a self-dual index choose a nonsingular invariant form of sign
$\varepsilon_i$, and compute the finite residue endomorphism field
$D_i$ with its adjoint involution $\sigma_i$. Define the symmetric,
alternating, unitary and non-self-dual-pair index sets
$\mathscr J_{\mathrm s},\mathscr J_{\mathrm a},
\mathscr J_{\mathrm u},\mathscr J_{\mathrm h}$ by the same explicit
conditions as in Theorem~\ref{C:Om-finite-field-algorithm}.
Then, with the concrete generators constructed below,
\[
 \begin{split}
 WG_{\mathrm{sp}}(BO_m/R_\nu)
   &\simeq
    \bigoplus_{i\in\mathscr J_{\mathrm s}(R_\nu)}
       \bigl(\mathbb Z e_i\oplus(\mathbb Z/2)t_i\bigr)\\
   &\hspace{8mm}\oplus
    \bigoplus_{i\in
      \mathscr J_{\mathrm a}(R_\nu)\cup
      \mathscr J_{\mathrm u}(R_\nu)\cup
      \mathscr J_{\mathrm h}(R_\nu)}
       \mathbb Z e_i,\\
 WG_{\mathrm{sp}}(BO_m/R)
   &\simeq \prod_{\nu=1}^s WG_{\mathrm{sp}}(BO_m/R_\nu).
 \end{split}                                             \tag{C.3.R1}
\]
The second line is an isomorphism of rings with exterior operations,
computed componentwise. Every requested product and every specified
exterior-power coefficient of a virtual element is obtained in
finite time. No uncomputed representation or form isomorphism
quotient remains in these indices or coefficients.
The indices over $R$ itself are obtained by extending each local
indecomposable by zero on the other factors.
\end{theorem}
\begin{proof}
\emph{Local decomposition of the base.}
Enumerate the ideals of the finite ring and identify its maximal
ideals by testing whether their quotients are fields. Their
intersection $J_R$ is nilpotent: its powers eventually stabilize,
and Nakayama applied to the stabilized finite module makes it zero.
Choose an exponent $a$ with $J_R^a=0$. Distinct maximal ideals are
comaximal, as are their $a$th powers, so the Chinese remainder theorem
gives
$R=\prod_{\mathfrak m}R/\mathfrak m^a$.
Each factor is local. All maps and inverse maps can be computed
from the finite tables. Finite projective modules, coactions, duals
and symmetric forms decompose over these base idempotents; cross
pairings between different components vanish. Tensor and exterior
powers are componentwise. This proves the last assertion in
(C.3.R1), once the local case is proved. The zero ring has zero
representation and form group and is handled separately.

\emph{Finite coefficient sets, without a free coalgebra assumption.}
Assume henceforth that $R$ is nonzero and local. Every finite
projective $R$-module is free. With the chosen split Gram matrix $J$ put
\[
 H=R[x_{ij}]/(X^{\mathsf t}JX-J).
\]
The usual comultiplication, counit and antipode are given by
$X\mapsto XY$, $X\mapsto I_m$, and
$X\mapsto J^{-1}X^{\mathsf t}J$, respectively. The determinant
relation makes $X$ invertible, so this is the full orthogonal
coordinate Hopf algebra.

Fix an order on the finite set $R$ and a monomial order. For each
$d$, enumerate all polynomials of total degree at most $d$; there
are finitely many. Use Lemma~\ref{C:finite-base-polynomial-equality}
to identify equal classes and retain the first representative in
the order which compares degree first and coefficient lists second.
Call the resulting finite set $F_dH$. Because degree is compared
first, these representative sets are nested. They represent the
image of the degree-at-most-$d$ polynomials, which is an $R$-submodule
of $H$. It is not asserted to be free. Their union is $H$, and
products, antipodes and all equality tests can be computed by the
same algorithm.

Enumerate all $n$-square matrices $h$ with entries in $F_dH$. Retain
exactly those satisfying
\[
 h(XY)=h(X)h(Y)
 \quad\text{in }M_n\bigl(R[X,Y]/(X^{\mathsf t}JX-J,\,
                                  Y^{\mathsf t}JY-J)\bigr),
 \qquad h(I_m)=I_n .                                    \tag{C.3.R2}
\]
These tests terminate by the equality lemma. They are precisely
coassociativity and counit for a coaction on $R^n$; the antipode
then supplies the inverse action matrix. The tensor-product
presentation in (C.3.R2) is valid without a flatness hypothesis.

Discard matrices already having all entries in $F_{d-1}H$ when
$d>0$. Polynomial coefficient degree is basis-independent, since
conjugation by a matrix over $R$ takes $R$-linear combinations
of the entries, and $F_{d-1}H$ represents an $R$-submodule.
Compute the centralizer by enumerating all $Z\in M_n(R)$ and
testing $Zh=hZ$ in $H$. Discard a coaction if this centralizer
has a nonzero nonidentity idempotent. The image and kernel of such
an idempotent are finite projective, hence free, invariant
summands. Conversely a decomposition gives such an idempotent.
Finally retain the least element of the finite
$\mathrm{GL}_n(R)$-conjugacy orbit, replacing entries by their
canonical polynomial representatives after each conjugation.
This constructs the claimed finite list without repetitions.
Every finite projective rational representation occurs because
its finitely many coefficients have a finite polynomial degree.

For actual computations, a commuting idempotent can be split
explicitly: enumerate $r$ and $P\in\mathrm{GL}_n(R)$ until
$P^{-1}ZP=\operatorname{diag}(I_r,0)$. The projective image and
kernel being free ensures success. Iterate on the two diagonal
coaction blocks; positive free ranks decrease strictly. This gives
a terminating decomposition into the canonical indecomposables.
Duality is computed by $h^\vee=\mathsf S(h)^{\mathsf t}$ and the
same finite canonicalization.

\emph{Why the hermitian-category theorem applies to a nonfield base.}
The additive category just used has split idempotents and finite
morphism sets. Every object is a finite set, so the Fitting
decomposition for an equivariant endomorphism terminates and
splits in this category. On an indecomposable, an endomorphism
is therefore either invertible or nilpotent. Its endomorphism
ring $E$ is local: if two nonunits summed to a unit, multiplying
by that inverse would give two nonunits summing to $1$, whereas
$1$ minus a nilpotent is a unit. Nonunits are also stable under
multiplication by arbitrary elements; a one-sided inverse in
a finite ring is two-sided. Thus they form the unique maximal
ideal. Repeated idempotent splitting gives a Krull--Schmidt
decomposition with these local endomorphism rings.

Every $E$ is finite, hence Artinian; its radical is nilpotent,
so it is complete for the radical topology. The dual of a finite
free module is finite free, biduality is exact on these objects,
and two is invertible. These verify all the hypotheses of
Quebbemann--Scharlau--Schulte
\cite[Propositions 2.4--2.5, p.271; Theorems 3.2--3.3,
pp.272--273]{C-QSS1979}, now in this additive category over $R$.
The hypothesis is not that the base itself be a field.

Compute each finite $E$ by the centralizer test. Its radical is
the set of matrices whose determinants are not units of $R$:
an invertible equivariant matrix has an equivariant inverse.
The quotient $D=E/\operatorname{rad}E$ is a finite division
ring, hence a finite field, as in
Theorem~\ref{C:Om-finite-field-algorithm}. Its characteristic is
odd because two is invertible. For a self-dual indecomposable,
enumerate all symmetric and alternating matrices over $R$,
test invariance by
$h^{\mathsf t}Bh=B$ in $H$, and test nonsingularity by
$\det B\in R^\times$. Proposition 2.5 guarantees success in at
least one sign. The adjoint involution on $E$ and on its quotient
field is computed from this chosen matrix.

Transfer and radical reduction identify the form contribution
with the corresponding finite-field multiplicity-form group.
The finite-field classifications and their group completions
were checked in Theorem~\ref{C:Om-finite-field-algorithm}:
symmetric type with trivial residue involution contributes
$\mathbb Z\oplus\mathbb Z/2$, alternating type contributes
$\mathbb Z$, nontrivial residue involution contributes $\mathbb Z$,
and a non-self-dual pair contributes a hyperbolic $\mathbb Z$.
This proves the additive formula in (C.3.R1).

The generators are concrete over $R$, with the same formulas
as over a finite field. If $q_i$ is symmetric with trivial residue
involution, let $e_i=[q_i]$. Lift a first nonsquare $c\in D_i$
to $b\in E_i$, put $a=(b+b^*)/2$, and set
$t_i=[q_i a]-[q_i]$. The reduction of $a$ is $c$, so $a$ is
invertible, and $q_i a$ is symmetric. In alternating or
non-self-dual-pair type take $e_i=H(U_i)$. In unitary type,
with chosen sign $\varepsilon_i$, choose
$0\ne c\in D_i$ with $\sigma_i(c)=\varepsilon_i c$,
lift it to $b$, and put $a=(b+\varepsilon_i b^*)/2$.
Then $a^*=\varepsilon_i a$, its reduction is $c$, and
$e_i=[q_i a]$ is an invertible symmetric form. Such $c$
exists, using $b-\sigma_i(b)$ for the negative sign.
All these searches are over explicitly finite sets.

\emph{Ring and exterior operations.}
For actual generators, multiply coaction entries for a tensor
product and take minors for an exterior power. These polynomials
have degree at most $d+e$ and $rd$, respectively. Canonical
polynomial equality and the preceding idempotent procedure
decompose the resulting actual module and identify all its
canonical indices and minimal degrees. The resulting actual
Gram matrices are tensor products and determinant forms, so
their computation uses no division by a factorial.

At the resulting fixed underlying multiplicities, radical
reduction leaves only a determinant bit for each nonempty
symmetric-type component and one forced class for the other
types. Construct the finitely many candidate orthogonal sums
of the displayed generators, exactly as in the finite-field
proof, and test each against the actual form by enumerating
all $P\in\mathrm{GL}_N(R)$ and checking equivariance and Gram
congruence. The hermitian-category theorem guarantees existence
and uniqueness of the successful residue-form coordinates.
This computes every structure constant in (C.3.R1), including
its two-torsion coordinates. Bilinearity and
$\lambda_t(x-y)=\lambda_t(x)\lambda_t(y)^{-1}$ give all virtual
operations to any prescribed finite degree. The equality
subroutines terminate, and every other enumeration is explicitly
finite. This completes the uniform algorithm.
\end{proof}

\begin{proposition}[Proved: a nonfield example over $\mathbb Z/9$]
\label{C:Om-finite-base-form-example}
Put $R=\mathbb Z/9$, let $V=R^3$ with Gram matrix $I_3$, and
let $S=\operatorname{Sym}^2V$ with its rational orthogonal action.
The group $O(I_3)$ is isomorphic to the split $O_3$ over $R$.
The module $S$ is indecomposable, and the
invariant symmetric forms
\[
                Q(A,B)=\operatorname{tr}(AB),\qquad 2Q
\]
give different classes whose difference has exact order two
in $WG_{\mathrm{sp}}(BO_3/R)$. Its isotypic contribution is
$\mathbb Z\oplus\mathbb Z/2$. More precisely,
\[
 \operatorname{End}_{O_3}(S)=R[I,K],\qquad
 K(A)=\operatorname{tr}(A)I_3,\qquad K^2=3K,\qquad
 E/\operatorname{rad}E=\mathbb F_3 .                     \tag{C.3.R3}
\]
\end{proposition}
\begin{proof}
The primitive vector $(1,1,4)$ is isotropic over $R$ and has
unit pairing with $(1,0,0)$. Subtracting half the norm times
the first vector gives a hyperbolic pair; the orthogonal
complement is a nonsingular line. Scaling the whole form by
the inverse of the norm on that line leaves its orthogonal
group unchanged and identifies it with the split odd-dimensional
orthogonal group.

The sign-change and permutation argument in
Proposition~\ref{C:Om-canonical-example} works over $R$ because
two is invertible. It constrains an endomorphism of symmetric
matrices to parameters $a,b,c$ there. To obtain $c=a-b$,
base change to the finite free, hence faithfully flat, algebra
$R'=R[z]/(z^2+1)$. Its elements $u=4z$, $v=z$ satisfy
$u^2+v^2=1$ and $uv=5\in R^\times$. The same plane-rotation
comparison therefore forces $c=a-b$ already over $R$.
The identity and $K$ are equivariant and $R$-linearly independent,
so this proves the presentation (C.3.R3).

The ideal $(3,K)$ has cube zero and the quotient is $\mathbb F_3$,
so $E$ is local. It has no nontrivial idempotents, proving
indecomposability. Both $I$ and $K$ are self-adjoint for $Q$.
The Gram matrix of $Q$ is
$\operatorname{diag}(1,1,1,2,2,2)$, hence invertible over $R$.
Every invariant nonsingular symmetric form on $S$ is therefore
$Q(aI+bK)$ with $a\in R^\times$, $b\in R$.
The residual involution is the identity, and the finite-field
transfer in Theorem~\ref{C:Om-finite-base-algorithm} gives
exactly the two rank-one multiplicity classes, represented by
$Q$ and $2Q$. Their residual coefficients $1,2\in\mathbb F_3$
are in different square classes. This also proves that their
difference cannot vanish after stabilization. The matrix
$\left(\begin{smallmatrix}1&1\\-1&1\end{smallmatrix}\right)$
has transpose-product $2I_2$ over $R$, so its tensor with $I_S$
gives $Q\perp Q\simeq2Q\perp2Q$. The difference has exact
order two.

The executable certificate
\texttt{code/B\_finish\_finite\_base\_forms.py} enumerates all
$81$ elements of $E$, all $54$ equivariant units and all $54$
nonsingular invariant symmetric forms. It checks that the only
idempotents are $0,I$, and its $2916$ isometry-action comparisons
give precisely two orbits of size $27$. It also verifies the
faithfully flat rotation identities and the rank-twelve
double-sum congruence. The complete symbolic endomorphism
argument above makes these enumerations exhaustive.
\end{proof}
